\clearpage
\SourceLeaf{174}\centerline{\Large BRIEFWECHSEL.}

\bigskip
\hrule
\bigskip

\centerline{[LEGENDRES THEORIE DER PARALLELEN.]}

\bigskip
\hrule
\bigskip

\centerline{[\textsc{Gerling} \emph{an} \textsc{Gauss}. \ \emph{Kassel, 11. März 1816.}]}

\smallskip

\{.\ .\ .\ .\ .\ .\ .\ Am Schlusse dieses fällt mir ein, dass ich schon oft vergessen habe, Sie um Ihr Urtheil über \textsc{Legendres} Theorie der Parallelen in seinen élémens de géom. zu bitten. --- Er definirt die gerade Linie als den kürzesten Weg zwischen zwei Punkten, und beweist mit Hülfe dieser Definition, dass die Summe 3er Winkel im Dreieck \textls{nicht grösser} sein kann als $2R$, nachher beweist er, dass sie auch \textls{nicht kleiner} sein könne als $2R$, wobei aber vorausgesetzt wird, dass man eine Linie durch einen Punkt zwischen den Schenkeln eines Winkels der $< \tfrac{2}{3}R$ ist, immer so legen könne, dass beide Schenkel geschnitten werden. Diese Voraussetzung rechtfertigt er in einer Anmerkung (pag. 280, 6te edit. 1806) durch das Einholen und Überholen des Punktes durch Verbindungslinien zwischen gleich weit vom Scheitel abliegenden Punkten auf den Schenkeln. --- In diesem letzten scheint mir aber derselbe
Fehler zu stecken, dessen Sie mich bei meinem letzten Aufenthalt in Göttingen überführten. Er erklärt es für \textls{assez évident}, und glaubt, man könne es zu keiner grössern Strenge bringen, ohne von einer andern Erklärung der geraden Linie auszugehen. --- Hinterher aber zeigt er, dass die Summe der 3 Winkel im Dreieck $= 2R$ sein müsse, noch auf eine andere mir neue und, wie mir scheint, stringente Weise; etwa so: durch 2 Winkel und die zwischenliegende Seite $A, B, c$ ist das ganze Dreieck bestimmt, also
\[
C = \varphi(A,B,c).
\]
